% ============================================================================
%  SPIE Methods section, subsections 2.1 - 2.7
%  Requires: \usepackage{amsmath, amssymb}   (for \mathcal, \sim, \times, \pm)
%
%  DRAFTING STATUS -- see SPIE_METHODS.md PENDING INVENTORY before submission.
%  Subsections 2.4 (correspondence probe) and 2.5 (student distillation) are
%  written in the past tense but HAVE NOT BEEN EXECUTED. The reported 84.7% /
%  75.5% are TEACHER numbers, measured with the policy reading privileged
%  state; if the student is trained, Results must be regenerated from it.
% ============================================================================

\section{METHODS}

\subsection{Procedural anatomy generation}

Anatomies are generated by perturbing a patient-derived arterial tree rather than by
sampling additional patient segmentations. The loaded tree comprises 25 centerline
branches; of these, only the navigated right common carotid artery (RCCA, total arclength
$237.5$\,mm) and the proximal $35$\,mm of the adjacent right vertebral artery (RVA) are
modified. All remaining branches are passed through unchanged.

\textbf{Displacement field.} For each generation, the RCCA centerline is displaced by a sum
of three sinusoidal modes in normalized arclength $u \in [0,1]$, with frequencies
$(0.7, 1.3, 2.1)$ cycles. Each mode draws an independent amplitude $a \sim \mathcal{N}(0,1)$
and phase $p \sim \mathcal{U}(0, 2\pi)$, and modes are summed independently along two
directions perpendicular to the branch's overall chord, then normalized by $\sqrt{3}$. The
displacement magnitude at a station is therefore itself a random quantity whose \emph{scale}
is set by a base amplitude of $4.0$\,mm multiplied by a per-generation tortuosity factor
$\tau \sim \mathrm{clip}\!\left(\mathcal{N}(1.0, 0.3),\, 0.4,\, 1.6\right)$.

\textbf{Envelope and anchoring.} The displacement is modulated by a bell envelope built from
two smoothstep ramps, $w = \min(w_{\mathrm{up}}, w_{\mathrm{down}})$, which is zero over the
first $3$\,mm and the last $25$\,mm of the branch and unity through the middle ($w = 1$ for
$s \in [18.0, 197.5]$\,mm, ramp length $15$\,mm). The proximal anchor pins the RCCA ostium,
which is shared with the RVA and with the parent bridge segment; the distal anchor pins the
terminus and its Circle-of-Willis junction. Because both ends are held fixed and the envelope
is $C^{1}$, the perturbed branch remains topologically connected to the unmodified tree and no
surface boolean operation is required. Lumen radii are scaled by the same envelope,
$\rho \leftarrow \rho\,(1 + (\sigma - 1)\,w)$ with
$\sigma \sim \max\!\left(0.75,\, \mathcal{N}(1.0, 0.07)\right)$, so calibre varies mid-vessel
while both junctions retain their original radius.

\textbf{Fork randomization.} The proximal $35$\,mm of the RVA receives an independent
perturbation (amplitude $3.0$\,mm, tapering to zero at $35$\,mm), so the geometry of the
RCCA/RVA bifurcation---and hence the deflection required to select the correct daughter
branch---differs between anatomies and cannot be memorized.

\textbf{Surface reconstruction.} The complete tree is re-meshed from centerlines at every
generation: spheres of the local radius are painted into a voxel grid of spacing
$[0.6, 0.6, 0.9]$\,mm, the volume is Gaussian-smoothed twice, an isosurface is extracted by
marching cubes, and the resulting mesh is decimated.

\textbf{Clearance audit.} Reconstruction narrows the lumen, so each generated surface is
audited before use. Clearance is measured as the distance from each centerline station to the
nearest point \emph{on the triangle surface}, sampled densely across faces; nearest-vertex
distance is not used, as it overstates clearance on a coarse mesh. An anatomy is admitted only
if the minimum clearance exceeds the guidewire radius ($0.18$\,mm) \emph{and} the median
clearance is at least $2.00$\,mm; anatomies failing either test are redrawn. A calibration
sweep established that a radius scaling of $1.6$ reproduces the patient-derived vessel's
clearance profile (median $2.14$\,mm versus $2.14$\,mm; fifth percentile $1.13$\,mm versus
$1.20$\,mm, i.e.\ marginally tighter). This ensures that anatomical variation is not obtained
by widening the vessel: a generated anatomy is no easier to traverse than the patient's own.

\textbf{Anatomy streams.} Each of 16 parallel simulation workers instantiates its own anatomy
and redraws it periodically, giving an estimated 700--800 distinct anatomies over a training
run.

\subsection{Simulation environment and control problem}

\textbf{Simulation.} Finite-element simulation uses SOFA with the BeamAdapter plugin under a
free-motion animation loop with Lagrangian constraint correction. The physics timestep is
$6$\,ms and 22 substeps are integrated per control step, giving a control rate of $7.5$\,Hz.
Contact uses a local-minimum-distance model (contact distance $0.3$\,mm, alarm distance
$0.5$\,mm) with a Coulomb friction coefficient of $0.001$.

\textbf{Devices.} A $0.36$\,mm J-tipped guidewire (tip radius $12.1$\,mm, tip arc $0.4\pi$\,rad)
and a $0.6/0.7$\,mm microcatheter, each $900$\,mm long and discretized into 553 beam elements,
are advanced coaxially. Both devices begin every episode at an independent uniform axial
rotation in $[0, 2\pi)$, which alone can determine which daughter branch the J-tip enters.

\textbf{Action.} Four continuous commands---translation rate ($\pm 30$\,mm/s) and twist rate
($\pm 1.5$\,rad/s) for each device---matching the degrees of freedom a bench actuator drives,
so that a learned policy is executable on hardware without re-parameterization.

\textbf{Observations, and the privileged/deployable split.} The full observation is
125-dimensional and partitions into a 101-dimensional \emph{deployable} prefix and a
24-dimensional \emph{privileged} block.

The deployable prefix contains quantities obtainable in a catheterization laboratory:
projected device tracking (40 dimensions; ten points at $15$\,mm spacing, tip-relative, with a
two-step history), target offset in the image plane (2), the previous command (4),
planned-route guidance (51), and the scripted controller's action for the current state (4).
Forty of the 51 guidance dimensions derive from the clinician's planned route: remaining
distance, cross-track offset, heading error, curvature ahead, distance and direction to the
next correct daughter branch, and an on-path indicator.

The privileged block contains simulator-only quantities: nodal velocities and forces at the
distal beam elements, a contact-impulse proxy formed from the difference between constrained
and free positions, accumulated axial twist, and ground-truth branch identity. These are
precisely the quantities that mark stress states---buckling, wall loading, and torsional
wind-up---and that are unavailable to a fluoroscopic observer.

\textbf{Reward.} A weighted sum of: a terminal target reward ($+3.0$ within $5$\,mm of the
target); dense arclength progress along the planned route ($0.01$ per mm, symmetric, so that a
round trip nets zero) with a lateral-deviation penalty ($0.001$ per mm); a small per-step cost
varying by path segment ($-0.007$ to $0$); discrete $\pm 1.0$ events on committing to a correct
or incorrect daughter branch at a fork; and a potential-based anti-buckling term (coefficient
$0.5$) formed from guidewire slack and mean contact displacement. Episodes terminate on target
contact and truncate at 600 control steps or on simulator error; vessel-end truncation is
disabled so that off-path excursions remain recoverable rather than fatal.

\textbf{Residual composition.} The executed command is
$a_{\mathrm{total}} = \mathrm{clip}\!\left(a_{\mathrm{heur}} + a_{\pi}\right)$ in raw physical
units. The scripted controller $a_{\mathrm{heur}}$ is a parameterless centerline follower: it
retracts when off-route, and otherwise corrects heading error and cross-track offset
proportionally while advancing, with the catheter following at $0.8\times$ the guidewire rate.
The residual has unity scale, so the network commands the full actuation range and can override
the controller entirely; at initialization, however, the system behaves exactly as the scripted
controller, which therefore serves as the study's baseline.

\subsection{Privileged teacher}

\textbf{Algorithm.} The teacher is trained with soft actor-critic using two critics and Polyak
targets ($\tau = 0.005$, $\gamma = 0.99$), taking $\min(Q_1, Q_2)$ in both the target and the
actor loss. The entropy term is omitted from the critic backup while retained in the actor
objective. Both networks are two-layer multilayer perceptrons of width 256 without recurrence;
layer normalization is applied in the critic bodies only. Optimization uses Adam at
$3 \times 10^{-4}$ with gradient-norm clipping at $1.0$. The entropy temperature is auto-tuned
against a target entropy of $+1.0$, rather than the conventional $-n_{\mathrm{actions}}$, with
$\log \alpha$ railed to $[-5, 0]$; the policy log-standard-deviation is soft-bounded to
$\sigma \in (0.135, 1.0)$.

\textbf{Choice of objective.} An advantage-weighted variant was used in an earlier generation of
this work and abandoned. Its exponential advantage weights collapsed toward unity, degenerating
the update into near-uniform behavior cloning of whatever the buffer contained---including the
policy's own increasingly extreme successes, a positive feedback loop with no counterweight.
Entropy-regularized soft actor-critic with a raised target entropy was adopted instead,
retaining an explicit exploration pressure that the advantage-weighted objective lacks. This
matters here because the recovery behavior characterized in Sec.~\ref{sec:behavior} is rare in
the buffer, and is precisely what an objective that clones the buffer's modal behavior would
discard.

\textbf{Replay.} A $2 \times 10^{6}$-transition buffer with prioritized sampling (exponent
$0.6$; importance-sampling exponent annealed $0.4 \rightarrow 1.0$) and batch size 256. Thirty
percent of each batch is drawn from a second tree restricted to episodes that reached the
correct daughter branch. There is no offline data, demonstration seeding, or pretraining: the
buffer is filled purely online after a $20{,}000$-step warm-up at reduced action scale.

\textbf{Scale.} Sixteen CPU workers feed one GPU trainer at a measured 10--11 environment steps
per second, with an update-to-data ratio of $0.99$. The teacher was trained for
$2.0 \times 10^{6}$ environment steps.

The teacher observes all 125 dimensions and is therefore not deployable; it exists to generate
supervision for the student of Sec.~\ref{sec:student}.

\subsection{Correspondence between planned-path geometry and privileged state}
\label{sec:correspondence}

% PENDING -- probe not yet implemented. See SPIE_METHODS.md inventory H1-H3.
% H2 (identity confound) is REAL: guidewire slack appears on both sides of the
% split, computed from the same projection. Excluding it is not optional.
% H3: existing informal evidence partly points the other way; this subsection
% may return a NEGATIVE result, which would require rewriting Sec. 2.5's motivation.

The distillation of Sec.~\ref{sec:student} is viable only if the deployable observation carries
information about the stress states on which the teacher conditions. We test this directly
rather than assume it.

A probe is trained to predict privileged stress quantities---contact-impulse magnitude and
distal nodal force---from the deployable planned-path features alone, on states sampled from
teacher rollouts across audited anatomies. Performance is reported as area under the receiver
operating characteristic curve for a binarized high-contact state, and as coefficient of
determination for the continuous targets, each against a shuffled-label control.

Two constraints make the test meaningful. First, any deployable feature that is numerically
identical to a privileged dimension is excluded from the probe: guidewire slack, in particular,
appears on both sides of the split and is computed from the same path projection, so including
it would measure an identity rather than a correspondence. Second, the probe is fitted and
evaluated on audited anatomies only, since geometry the device cannot traverse produces contact
signatures that no deployable feature should be expected to explain.

The result establishes the extent to which planned-path geometry acts as an observable proxy for
unobservable mechanical state, and therefore bounds what distillation can recover.

\subsection{Student distillation}
\label{sec:student}

% PENDING -- no distillation code, launcher, run directory or checkpoint exists.
% See SPIE_METHODS.md inventory S1-S4 and R1.
% S4: building the student requires changing privileged_obs_dim from 0 -- the same
% flag that silently made the "asymmetric" actor symmetric in the teacher runs.
% R1: 84.7% / 75.5% are TEACHER numbers. If this subsection ships, Results MUST be
% regenerated from the student.

The deployable controller is obtained by distilling the teacher into a student restricted to the
101-dimensional deployable prefix, using dataset aggregation. The student is initialized by
regression to the teacher's actions on teacher-generated states, and then updated iteratively:
the student is rolled out in the environment; the teacher, which retains access to privileged
state at every visited configuration, labels each visited state with the action it would have
taken; and the student is updated by regression toward those labels on the aggregated dataset.
Iterating on states the \emph{student} visits, rather than only those the teacher visits,
corrects the compounding distribution shift that pure behavior cloning suffers when the
student's errors carry it into configurations the teacher never entered.

The student shares the teacher's architecture apart from the input width, and is composed with
the same scripted controller under the same residual rule, so that teacher and student differ
only in what they observe. All reported deployable results are the student's; the teacher's
score is reported alongside as an upper bound on what distillation could recover.

\subsection{Behavioral analysis}
\label{sec:behavior}

% PENDING re-analysis -- the correlational result motivating this framing was
% computed on the pre-audit anatomy distribution, in which a large fraction of
% anatomies were geometrically impassable. See SPIE_METHODS.md inventory B1.

Navigation performance is decomposed into two behaviors by post-hoc analysis of the complete
step logs. A \emph{stall} opens when the guidewire is advanced above $2$\,mm/s while path
progress remains flat within $0.3$\,mm for twelve consecutive steps, using a decaying counter so
that isolated command reversals do not close it. A stall resolves when progress passes the
frontier at which it began, and is classified by the guidewire length actually withdrawn---under
$1$\,mm, $1$--$8$\,mm, or over $8$\,mm---measured from executed insertion rather than commanded
retraction, since a retraction command against a wedged device moves nothing.

Three success measures are separated because they occupy different units of analysis: episode
success (per episode), stall-resolution rate (per stall event), and the conversion of resolved
stalls into successful episodes (per episode). Absolute rates move by roughly 20 percentage
points across reasonable threshold choices, so only differences and trends are reported, and the
analysis is restricted to audited anatomies.

\subsection{Evaluation protocol}
\label{sec:eval}

Evaluation uses 98 fixed seeds across 50 procedurally generated anatomies, redrawn every two
episodes across 16 workers, with the clearance audit of Sec.~2.1 applied at generation time.
Actions are taken deterministically as the hyperbolic tangent of the policy mean. An episode
succeeds if the guidewire tip comes within $5$\,mm (three-dimensional Euclidean distance) of a
target sampled at least $40$\,mm into the RCCA, within 600 control steps.

Anatomy identity is established by hashing branch-coordinate arrays. This is necessary because
the obvious identifiers are unreliable: the planned path moves with the target, and the mesh
fingerprint is built from a seed that is reassigned on reset, so either can report a match
between geometrically different vessels.

Results are stratified by planned-path length into three bands, with cut points at $146$\,mm and
$210$\,mm, corresponding approximately to the common carotid artery, the mid-internal carotid
artery, and the cavernous segment. These are cut points on path length rather than anatomical
segmentation; the modelled vessel terminates at the internal carotid artery terminus, and the
middle cerebral artery is not represented.

Transfer is assessed on the patient-derived anatomy using its original segmented surface rather
than a reconstruction, with all other settings unchanged.

Binomial (Wilson) confidence intervals are reported. These understate the true uncertainty:
episodes are not independent, as 16 persistent simulator processes each run six to seven
sequential episodes, and the initial device twist is unseeded, so repeat launches are not
bit-reproducible.
